\section{未来研究展望}

第 5 节将路线比较收束为若干系统挑战，本节进一步讨论未来更值得持续投入的研究方向。这些方向可压缩为三条主线：面向异构平台的能力契约与共享抽象、面向真实交付的控制面与证据链治理，以及面向复杂请求图的状态中心化运行时。图\ref{fig:outlook-map} 概括了三条主线在国产推理生态中的收敛关系。

\begin{figure}[!htb]
\centering
\resizebox{0.94\linewidth}{!}{\begin{tikzpicture}[
    node distance=18mm and 24mm,
    block/.style={draw=black!75, line width=0.8pt, rounded corners=2.2pt, minimum width=31mm, minimum height=10mm, align=center, font=\small, inner sep=4pt},
    centerblock/.style={draw=black!75, line width=0.8pt, rounded corners=3pt, minimum width=40mm, minimum height=12mm, align=center, font=\small\bfseries, inner sep=5pt},
    note/.style={font=\footnotesize, align=center, fill=white, inner sep=1.2pt, text=black!80},
    flow/.style={-{Latex[length=2.2mm]}, semithick, draw=black!75}
]
\node[centerblock, fill=gray!8] (goal) {可持续的国产算力\\推理系统底座};

\node[block, fill=blue!9, above left=20mm and 28mm of goal] (plugin) {后端插件化\\平台抽象/注册/可合并优化};
\node[block, fill=green!10, above right=20mm and 28mm of goal] (runtime) {资源管理与调度\\SLO/迁移/KV Cache/阶段解耦};
\node[block, fill=orange!12, below left=20mm and 28mm of goal] (eval) {统一评测与观测\\近真实负载/成本/稳定性};
\node[block, fill=red!8, below right=18mm and 24mm of goal] (model) {模型结构持续演进\\MoE/MLA/多模态/语音交互};

\draw[flow] (plugin.south east) -- (goal.north west);
\draw[flow] (runtime.south west) -- (goal.north east);
\draw[flow] (eval.north east) -- (goal.south west);
\draw[flow] (model.north west) -- (goal.south east);

\node[note, text width=6.2cm, above=6mm of goal] {未来竞争将从单点热点优化转向系统能力收敛。};
\node[note, text width=6.4cm, below=6mm of goal] {平台抽象、运行时、评测与模型演进需要同步推进。};
\end{tikzpicture}}
\caption{国产推理引擎未来演进方向关系图}
\label{fig:outlook-map}
\end{figure}

这三条主线并不是彼此独立的新功能列表，而是前文"一个底座、三条主路径"在未来更可能沉淀出的高层研究议题。它们之间存在递进关系：能力契约为系统演进提供稳定边界，控制面治理让这些边界可验证、可交付，而请求图重构则要求前两者同时升级以适应更复杂的运行时状态。

\subsection{面向异构平台的能力契约与共享抽象}

核心问题是：平台差异能否从"项目内部经验"上升为"生态可以共享的能力契约"。具体而言，需要回答以下开放研究问题：

\begin{enumerate}[leftmargin=*, label=(\arabic*)]
\item 后端能力应如何形式化表达，才能既覆盖 graph/eager 约束、通信语义和低比特路径的真实差异，又不把短期平台特例固化为长期接口？现有工程实践中，能力边界通常以隐式特判存在，缺少统一描述语言。
\item 回退语义应如何定义，才能保证图模式切换、缓存状态和分布式执行在回退前后保持一致？当前系统对回退的处理通常是"降级到 eager"，缺乏系统性的状态保持保证。
\item 共享抽象如何与平台特化协同演进？若抽象过薄，平台差异反复溢出；若过厚，又把短期特例固化。Llumnix、SOLA、LServe 等工作之所以重要，正因为它们显示运行时组织方式正在脱离"单一模型 serving 外壳"，转向资源与状态中枢的形态\citep{sun2024llumnix,hong2025sola,yang2025lserve}。
\end{enumerate}

只有当平台差异能够被收敛为共享能力矩阵而非反复改写系统主链时，新模型、新后端和新任务形态的进入才会从高风险迁移变成稳定的生态扩展。然而，能力契约本身并不能自动保证这些边界在生产环境中被持续遵守——这需要控制面来承担验证、回归和交付的闭环责任。

\subsection{控制面、评测与交付的一体化治理}

核心问题是：控制面能否从运维脚本集合演化为连接评测、成本与交付的一体化治理层。具体开放问题包括：

\begin{enumerate}[leftmargin=*, label=(\arabic*)]
\item 近真实工作负载应如何标准化表达？当前缺少能同时覆盖请求分布、上下文长度变化、多模态混合和工具调用比例的统一 workload 描述格式。BurstGPT 等工作已开始提供真实负载数据集\citep{wang2025burstgpt}，但从数据集到可复现评测协议之间仍有显著距离。
\item 版本基线、回退语义与故障注入应如何统一？升级回归和灰度成本目前无法被量化治理，大多停留在人工经验阶段。
\item 模型效果、系统成本与交付代价如何进入同一证据链？避免"性能结论"与"生产结论"长期脱节。
\end{enumerate}

执行面负责 token 生成，控制面负责让 token 生成处于可比较、可复现、可灰度和可回滚的条件之下。只有控制面被视为与执行面并列的一等对象，前文讨论的成本、评测与交付问题才可能从局部经验转化为稳定的生态能力。控制面治理的成熟度，反过来也取决于运行时本身能否提供更清晰的状态对象——这正是第三条主线关注的核心问题。

\subsection{面向多阶段请求图的运行时重构}

核心问题是：随着文本 decode、视觉 encode、语音前端、结构化约束验证、工具调用回注和多 Agent 协同被纳入同一服务链路，推理引擎如何从维护单一 token 生成循环转向管理多阶段请求图。具体开放问题包括：

\begin{enumerate}[leftmargin=*, label=(\arabic*)]
\item 请求图中的状态转移应如何显式建模？DeepSeek-V2/V3 的 MLA 与 MoE 组织连同 FlashMLA、Mooncake 等配套工作，已把通信组织、缓存布局和注意力实现重新推回运行时中心位置\citep{liu2024deepseekv2,deepseekv3,deepseekflashmla,qin2024mooncake}。这要求状态对象不再只按模型名单组织，而应围绕阶段关系组织。
\item 共享缓存与编译产物应如何围绕阶段关系而非模型名单组织？新任务形态进入时不应需要重写主链。
\item 运行时如何同时容纳模型 co-design、状态治理和交付控制面？Qwen2-VL、InternVL2、GLM-4-Voice 与 MiniCPM-o 等模型已说明视觉、视频和语音链路使文本时代的单一路径假设难以维持\citep{wang2024qwen2vl,chen2024internvl2,zeng2024glm4voice,cui2026minicpmo45}。
\end{enumerate}

这个转变未必立刻带来最亮眼的跑分，但更可能决定国产推理系统能否在未来几年内持续吸收新结构、新任务和新平台。系统边界会被模型公司、应用形态和硬件路线共同牵引，任何一方的变化都会迫使另外两方重写分工。

综合来看，这三条主线共同指向一个判断：国产推理引擎的下一阶段竞争，将从"谁更快适配某个硬件后端"转向"谁能在异构平台、复杂负载和快速演进之间建立更稳定的工程秩序"。能力契约、控制面治理和请求图重构不是三个孤立议题，而是同一问题的三个切面——系统能否把持续增加的复杂度组织进可扩展的边界，而不是让它们在外围不断堆积。
